Estas son algunas preguntas que surgen de la práctica, con su respuesta.

\textbf{1. ¿Por qué el rayo no se dobla al salir por la cara curva de la lente?}

Porque dentro de la lente viaja a lo largo de un radio del semicírculo, y un radio siempre llega perpendicular a la curva. Si el ángulo de incidencia es cero, la ecuación \eqref{eq:snell} da un ángulo de refracción también de cero.

\textbf{2. ¿Qué pasa con un ángulo de incidencia de 0 grados?}

El rayo atraviesa la lente sin desviarse. Por eso se usó esa posición para alinear el equipo antes de empezar.

\textbf{3. ¿Qué ocurriría si la luz entrara primero por la cara curva?}

Entonces iría del acrílico al aire en la cara plana, se alejaría de la normal y, al pasar el ángulo crítico de unos 42 grados, ya no saldría: se reflejaría por completo dentro de la lente. Ese fenómeno, la reflexión total interna, es el que permite que funcione la fibra óptica.

\textbf{4. ¿El índice de refracción depende del color de la luz?}

Sí, un poco. La luz violeta se dobla más que la roja en casi todos los materiales, y por eso un prisma separa la luz blanca en un arcoíris. Con una fuente de luz blanca como la de esta práctica, el rayo refractado puede verse ligeramente coloreado en los bordes.
